\section{What a Shared Prototype Operator Can Do to a Noisy Vote}
\label{sec:operators}

% Pipeline schematic (D-146): where the adversary, the noise and each registered operator family act. No data.
\begin{figure*}[t]
\centering
\begin{tikzpicture}[
  font=\scriptsize,
  box/.style={draw=black!70, rounded corners=2pt, align=center, inner sep=2.4pt, minimum height=11.5mm, fill=white},
  op/.style={box, fill=orange!12, draw=orange!75!black},
  sup/.style={box, fill=blue!6, draw=blue!60!black, dashed},
  cert/.style={box, fill=green!8, draw=green!45!black},
  arr/.style={-{Stealth[length=3.2pt]}, semithick},
  note/.style={align=center, font=\tiny, text=black!70},
  node distance=3.2mm,
]
% row 1: the certified path of one input
\node[box] (in) {model input\\ $x\in[0,1]^{3\times224\times224}$\\ adversary: $x+\delta$, $\lVert\delta\rVert_2<r$};
\node[box, right=of in] (noise) {Gaussian noise\\ $\varepsilon_i\sim\mathcal N(0,\sigma^2 I)$\\ draws $i=1,\dots,n$};
\node[box, right=of noise] (enc) {CLIP image\\ encoder\\ feature $z_i$};
\node[op, right=of enc] (imgop) {image side\\ $z'_i=\Normalize(z_i-c\,\mu_I)$\\ {\tiny (image half)}};
\node[box, right=of imgop] (vote) {hard-label vote\\ of draw $i$\\ $\argmax_k\langle z'_i,t'_k\rangle$};
\node[cert, right=of vote] (cert) {counts over $n$ draws\\ Clopper--Pearson $p_L$\\ $R=\sigma\,\Phi^{-1}(p_L)$, or abstain};
\foreach \a/\b in {in/noise, noise/enc, enc/imgop, imgop/vote, vote/cert} {\draw[arr] (\a) -- (\b);}
% row 2: prototypes and the operators that replace them
\coordinate (row2) at ($(in.south)+(0,-7mm)$);
\node[box, anchor=north] (txt) at (row2 -| enc) {class prompts\\ CLIP text encoder\\ prototypes $t_k$};
\node[op, anchor=north west, right=of txt.north east, anchor=north west] (textop) {text side, one common $w$ for all classes\\
  $t'_k=\Normalize(t_k+w)$\\
  {\tiny mean-centering $w=-c(\mu_T-\mu_I)$; text half $w=-c\,\mu_T$}\\
  {\tiny projected gap $w=c\,p$; boundary-active and noisy-margin $w=\eta u$}};
\node[sup, anchor=north west, right=of textop.north east, anchor=north west] (sup) {supervised controls\\ (use labels): rank-8 adapter,\\ class means, few-shot prompts,\\ learned shared $w$};
\draw[arr] (txt.east |- textop.west) -- (textop.west);
\draw[arr] (textop.north) -- ++(0,2mm) -| ($(vote.south)+(-2mm,0)$);
\draw[arr, dashed, blue!60!black] (sup.north) -- ++(0,3.4mm) -| ($(vote.south)+(2mm,0)$);
\node[note, above=0.8mm of noise] {noise and $\delta$ act before channel standardization};
\node[note, above=0.8mm of imgop] {two-sided centering = text half + image half};
\end{tikzpicture}
\caption{Where each registered operator acts. The adversary's perturbation
$\delta$ and the Gaussian noise live in the $224\times224$ model-input space.
Every shared text-side operator translates all prototypes by one common vector $w$
before renormalization, so by Proposition~\ref{prop:flip} it can change a
draw's vote only through the renormalization factors; image-side centering
changes the feature itself. The supervised controls (dashed) use labels; the
class-specific ones bound what is recoverable outside the shared family.}
\label{fig:pipeline}
\end{figure*}


\subsection{The registered operator family}
\label{sec:family}

All operators are fixed per cell from development data only and then frozen;
Figure~\ref{fig:pipeline} shows where each acts. With $\Normalize(v)=v/\lVert
v\rVert$, a coefficient $c\in\{0.25,0.5,1\}$ or a step
$\eta\in\{0.0025,0.005,0.01,0.02,0.04\}$, they are:
\begin{itemize}
\item \emph{Text mean-centering}: $t'_k=\Normalize\big(t_k-c\,(\mu_T-\mu_I)\big)$.
\item \emph{Projected-gap translation} (Chowers et al.~\cite{chowers2026bug}):
$t'_k=\Normalize(t_k+c\,p)$, where $p=(I-VV^{\top})(\mu_I-\mu_T)$ is the
component of $\mu_I-\mu_T$ orthogonal to all differences $t_i-t_j$
($V$ an orthonormal basis of their span).
\item \emph{GR-CLIP-style two-sided centering}: $t'_k=\Normalize(t_k-c\,\mu_T)$
together with $z'=\Normalize(z-c\,\mu_I)$ for every clean or noisy image
feature; the audit used $c=1$, and the follow-up also evaluates its text half
(prototypes only) and its image half ($z$ only).
\item \emph{Clean boundary-active translation}: $t'_k=\Normalize(t_k+\eta u)$
with a unit direction $u$ that increases the normalized clean margin of the
development items' own raw predictions.
\item \emph{Noisy-margin translation}: the same form with a unit direction fitted
to increase a soft noisy-margin objective over Gaussian draws of the development
items at the audit noise level.
\end{itemize}
Together with the identity this gives the \ExpCandidates{} registered banks per
cell. Ground-truth labels are never used to build a direction; the two proposal
directions target each item's own raw clean prediction.

\subsection{Per-draw identity and a decision-change condition}
\label{sec:flip}

Every text-side operator above has the form $t'_k=(t_k+w)/n_k$ with one common
translation $w$ and per-prototype renormalization factors
$n_k=\lVert t_k+w\rVert$. We derived
them after the audit outcome was known, so they explain the audit in hindsight
rather than predicting it; a weaker form of Proposition~\ref{prop:flip} was
registered as a conformance prediction for the extension of
Section~\ref{sec:extension} before its sampling.

% sections/theory_decision_change.tex
% Input from sections/operators.tex inside the subsection labeled sec:flip; no sectioning command here.
% Environments proposition/corollary/remark and the macros \Normalize, \argmax are defined in main.tex.
% Full proofs: sections/appendix_proofs.tex (the Proofs appendix, label app:proofs).
% Derivations, exact code references and the numerical tests behind every statement:
% research/DECISION_CHANGE_CONDITION.md. No result value is typed in this file.

Throughout, $z\in\mathbb S^{d-1}$ is one image feature (the clean feature or
one normalized noisy draw), $s_k=\langle z,t_k\rangle$ are its original scores,
$\omega=\langle z,w\rangle$, and every $n_k$ is positive.

\begin{proposition}[Per-draw identity]
\label{prop:identity}
For every $z$, $\langle z,t'_k\rangle=(s_k+\omega)/n_k$. Hence for every pair
$(j,k)$
\begin{align}
\langle z,t'_j\rangle-\langle z,t'_k\rangle
&=\frac{n_k(s_j+\omega)-n_j(s_k+\omega)}{n_j\,n_k},
\label{eq:identity}\\
n_j-n_k&=\frac{2\,\langle t_j-t_k,\,w\rangle}{n_j+n_k}.
\label{eq:normalizer}
\end{align}
\end{proposition}
\noindent\emph{Sketch.} Expand $\langle z,t_k+w\rangle$, and
$n_j^2-n_k^2=\lVert t_j+w\rVert^2-\lVert t_k+w\rVert^2=2\langle t_j-t_k,w\rangle$
because $\lVert t_j\rVert=\lVert t_k\rVert=1$.

\begin{proposition}[Necessary condition for a decision change]
\label{prop:flip}
Let $s_j\ge s_k$ and suppose that after the translation
$\langle z,t'_k\rangle\ge\langle z,t'_j\rangle$. Then
\begin{align}
0\le s_j-s_k
&\le (n_j-n_k)\,\langle z,t'_j\rangle
\le\lvert n_j-n_k\rvert
\nonumber\\
&=\frac{2\,\lvert\langle t_j-t_k,w\rangle\rvert}{n_j+n_k}
\le\min\{1,\lVert w\rVert\}\,\lVert t_j-t_k\rVert .
\label{eq:flip}
\end{align}
Consequently, for $\lVert w\rVert=\eta$ any pair whose order changes has
original cosine margin at most $\eta\lVert t_j-t_k\rVert\le 2\eta$, and if
the top-1 label of $z$ changes from $j$ to $k$, then the original top-1 margin
of $z$ (winner minus runner-up) is at most
$\lvert n_j-n_k\rvert\le\max_{i,l}\lvert n_i-n_l\rvert$.
\end{proposition}
\noindent\emph{Sketch.} Rearranging the hypothesis with
\eqref{eq:identity}--\eqref{eq:normalizer} gives the first two inequalities,
$\lvert\langle z,t'_j\rangle\rvert\le1$ gives the third, and the last one is
Ptolemy's inequality for the points $t_j$, $t_k$, $-w$, $0$ together with the
reverse triangle inequality (Appendix~\ref{app:proofs}). The second inequality
also carries a direction: when $\langle z,t'_j\rangle>0$ and $s_j>s_k$, the original winner
can lose only to a class whose translated raw norm is strictly smaller,
$n_k<n_j$.

\begin{corollary}[Projected-gap translation is decision-invariant]
\label{cor:chowers}
If $\langle w,t_i-t_j\rangle=0$ for all $i,j$, as for the projected-gap
translation $w=c\,p$ at every coefficient $c$, then all $n_k$ equal one value
$n$ and $\langle z,t'_k\rangle=(s_k+\omega)/n$ with $n$ and $\omega$ independent of $k$,
so the $\argmax$ (including ties) is unchanged for every $z$. In exact
arithmetic with a fixed tie rule, on any fixed set of noisy draws the operator
therefore reproduces the uncorrected hard labels,
vote counts, selected classes, certified radii, and accuracies exactly: it is
an exact invariance control, not a competing method.\footnote{This is the
mechanism behind the identical aggregates of the three projected-gap
coefficients and of the uncorrected bank in the saved audit record; in
floating-point arithmetic only numerical ties could differ.}
\end{corollary}
\noindent\emph{Sketch.} $n_k^2=1+2\langle t_k,w\rangle+\lVert w\rVert^2$ and
$\langle t_k,w\rangle=\langle t_1,w\rangle$ for every $k$.

\begin{corollary}[Finite flip budget]
\label{cor:budget}
Fix $N$ draws $z_1,\dots,z_N$, let $v_i$ be the number of draws whose original
$\argmax$ is $i$, let $k_0=\max_i v_i$, and let $E_F$ be the number of draws
whose original winner $j$ has a competitor $k\ne j$ with
$s_j-s_k\le\lvert n_j-n_k\rvert$; $E_F$ is at most the number of draws whose
original top-1 margin is at most $\max_{i,l}\lvert n_i-n_l\rvert$. After the
translation the counts $v'_i$ satisfy $v'_i\le v_i+E_F$ for every $i$, hence
$\max_i v'_i\le k_0+E_F$. In particular, if $k_0+E_F<k_{\min}$ for a
confirmation threshold $k_{\min}$, then no class reaches $k_{\min}$ on these
draws after the translation.
\end{corollary}
\noindent\emph{Sketch.} By Proposition~\ref{prop:flip} every flipped draw is
one of the $E_F$ eligible draws, and a class gains at most one vote per flipped
draw. Because $E_F$ depends only on the original scores and on the normalizers
$n_k$ recorded when a bank is built, the corollary can rule out a threshold
crossing on these draws, without rescoring, whenever $k_0+E_F<k_{\min}$; when
the inequality fails, it does not show that a crossing occurs. Saved top-1
margins give a looser upper bound on $E_F$, and the exact pairwise count needs
the original score differences. The certificate itself is still computed only
from the hard-label counts by the separate Cohen procedure; the bound is not a
certificate.

\begin{remark}[Two-sided operators]
\label{rem:twosided}
The GR-CLIP-style operator also maps $z\mapsto z'=\Normalize(z-\mu_I)$; we write the
audit's coefficient $c=1$, and a general $c$ replaces $\mu_I$ and $\mu_T$ by $c\mu_I$ and $c\mu_T$.
Renormalizing $z$ never changes an $\argmax$, but the shift does. Assume
$z\neq\mu_I$ and $t_k\neq\mu_T$ for every class $k$, so that both
normalizations are defined. With $n_k=\lVert t_k-\mu_T\rVert>0$,
\begin{align}
\langle z',t'_k\rangle
&=\frac{\langle z-\mu_I,\,t_k-\mu_T\rangle}{\lVert z-\mu_I\rVert\,n_k}
\nonumber\\
&\propto
\frac{s_k-\langle z,\mu_T\rangle-\langle\mu_I,t_k\rangle+\langle\mu_I,\mu_T\rangle}{n_k},
\label{eq:twosided}
\end{align}
where the proportionality constant is common to all classes. The
class-dependent term $-\langle\mu_I,t_k\rangle$ is a per-class bias that no
common text translation produces, so
Propositions~\ref{prop:identity}--\ref{prop:flip} do not bound this operator:
it can change the label of a draw whose original margin exceeds every bound in
\eqref{eq:flip}. This is why, within the audit's registered family, it is the only operator
that can act where the text-only translations cannot.
\end{remark}

\begin{remark}[Engine boundary]
Nothing in this section is a certificate; certified radii are computed
separately from hard-label counts as in Section~\ref{sec:background}.
\end{remark}


Two consequences matter for reading the results. A step of $\eta=0.0025$ can
flip a noisy draw only if its original top-1 margin is at most
$\lvert n_j-n_k\rvert\le2\eta$. The bound defines which draws are eligible;
how many are eligible depends on the empirical margin distribution, and on the
saved draws the registered steps changed very few certified outcomes
(Section~\ref{sec:nearidentity}). The projected-gap
translation leaves every hard label unchanged for every coefficient, so under
hard-label smoothing it is an exact negative control rather than a competing
method; this is a property of the projection Chowers et al.\ designed to
preserve nearest neighbors, made explicit for the smoothed classifier. Within
the audit's registered family, only the two-sided operator, which also shifts
the image feature, escapes the bound.

\subsection{Why a clean global-gap summary cannot settle the question}
\label{sec:nonident}

One might hope to predict a certificate from the clean geometry alone. It cannot
be done in general. Appendix~\ref{app:nonident} gives an explicit smooth
normalized encoder and two binary prototype banks that share the encoder, the
clean feature of the certified input, the complete global-gap vector under any
common image-reference distribution, the clean margin and the prototype
correlation, yet have population Cohen radii that differ by an arbitrary factor.
The two banks are mirror images of one another, and the only asymmetry is how
the encoder's response direction meets each decision boundary; the listed clean
summaries do not contain that response. This is a scoping remark, not a theorem about
CLIP: it motivates measuring noisy votes rather than clean gaps.
